\documentclass[11pt]{book}

\usepackage{graphicx}
\usepackage{eso-pic}
\usepackage{tcolorbox}
\usepackage{polyglossia}
\setdefaultlanguage[indentfirst=false,forceheadingpunctuation=false]{russian}
\setotherlanguages{english}

\RequirePackage[a4paper, headsep = 0.5 \headsep, left=2.5cm, right=2.1cm, top=2cm, bottom=2.1cm]{geometry}

\setmainfont{Times New Roman}
\newfontfamily\cyrillicfont{Times New Roman}[Script=Cyrillic, Ligatures=TeX]

\usepackage{enumitem}
\setlist[itemize, 1]{label =\raisebox{-0.3\height}{\scalebox{1.8}{\textbullet}}}

\title{Феноменология: Основы}
\author{Дэн Захави}
\date{2019}
\begin{document}

\AddToShipoutPictureBG*{\includegraphics[width=\paperwidth]{title.png}}

\null
    \thispagestyle{empty}%
    \addtocounter{page}{-1}%
    \newpage

\tableofcontents

\chapter{Введение}

Феноменология считается одной из доминирующих традиций в философии XX века. Ее основал Эдмунд Гуссерль (1859-1938)\footnote{Эдмунд Гуссерль пришел в философию в относительно позднем возрасте. Он родился в Простеёве, Моравия (тогда входившая в состав Австрийской Империи), изучал физику, математику, астрономию и философию в Лейпциге, Берлине и Вене. В 1883г он защитил докторскую диссертацию по математике, и в последующие годы серьезно заинтересовался философией , отчасти благодаря посещению лекций психолога и философа Франца Брентано. Первая крупная работа Гуссерля --- ''Логические исследования'' --- была опубликована в 1900-1901гг, и именно на ее основе он получил приглашение в Геттинген, где преподавал с 1901 по 1916 год. Следующая крупная работа --- ''Идеи к чистой феноменологии и феноменологической философии I'' --- была опубликована в 1913г. В 1916г Гуссерль переезжает во Фрайбург, где получает кафедру философии, прежде занимаемую Генрихом Риккертом. В последующие годы его ассистентами работают Эдит Штайн и Мартин Хайдеггер, который, когда Гуссерль выйдет на пенсию в 1928г, займет его место. Вскоре после выхода на пенсию Гуссерль публикует ''Формальную и трансцендентальную логику'' (1929) и ''Картезианские размышления'' (1931). В тридцатые годы Гуссерль страдал от условий, навязанных немецким национал-социалистическим режимом: лишенный какой-либо официальной академической деятельности из-за своего еврейского происхождения, он теряет право преподавать, публиковаться, и, в конечном итоге, свое немецкое гражданство. Он, однако, продолжает свою работу, в такие темные времена еще более настаивая на актуальности философии. В 1935г Гуссерля приглашают читать лекции в Вене и Праге, эти лекции в последствии лягут в основу его последней работы ''Кризис европейских наук и трансцендентальная феноменология'' (1936). Вскоре после его смерти 27 апреля 1938г молодому монаху-францисканцу Герману Лео ван Бреде удается контрабандой вывезти многочисленные исследовательские рукописи Гуссерля из Германии в монастырь в Бельгии. Перед началом Второй Мировой Войны в Левене был создан Архив Гуссерля, где до сих пор можно найти оригинальные рукописи.}, другие влиятельные сторонники --- Мартин Хайдеггер (1889-1976)\footnote{Хайдеггер начал изучать католическую теологию и средневековую философию во Фрайбурге, но в 1911г решил сосредоточиться на философии. В 1913г он защищает диссертацию, а два года спустя --- хабилитацию ''Учение Дунса Скотта о категориях и значении'', которая затем будет передана Риккерту --- философу, преемником которого был Гуссерль. Хайдеггер работает ассистентом Гуссерля в 1918-1923гг, и затем становится профессором университета в Марбурге. В 1927 году он публикует главную свою работу --- ''Бытие и время'' --- а в 1928 занимает кафедру Гуссерля во Фрайбурге. В 1929г Хайдеггер прочел свою знаменитую вступительную лекцию ''Что такое метафизика?''. После прихода к власти Гитлера Хайдеггер был избран ректором Фрайбургского университета и стал членом нацистской партии. Однако менее чем через год он уходит с этого поста и постепенно отходит от университетской политики. Как показывает публикация в 2014г его так называемых ''Черных тетрадей'', флирт Хайдеггера с нацизмом не был ни кратковременным, ни поверхностным. До 1944г он регулярно читал лекции, но после окончания войны ему запретили преподавать из-за политических убеждений, а в 1946г он был лишен звания профессора. В 1949г его восстанавливают в звании почетного профессора, и с тех пор и до своей смерти о читает много лекций и публикует такие центральные работы как ''Язык'' (1950), ''Размышления о строительстве и обитании'' (1951) и ''Вопрос технологии'' (1953).}, Жан-Поль Сартр (1905-1980)\footnote{Сартр изучал философию в Высшей нормальной школе, где встретил целое поколение ведущих французских интеллектуалов, в том числе Симону де Бовуар, Раймона Арона, Мориса Мерло-Понти, Симону Вейль, Эммануэля Мунье, Жана Ипполита и Клода Леви-Стросса. Его отношения с коллегой-философом Бовуар станут легендарными. В период 1931-1945гг Сартр преподает в средних школах Гавра, Лана и Парижа. С философией Гуссерля и Хайдеггера его знакомят в начале 1930гг Арон и Левинас, а в 1933-1934гг он отправляется в Берлин изучать феноменологию. Во второй половине тридцатых годов Сартр публикует четыре книги о сознании: ''Трансцендентность эго'', в которой рассматривается структура самосознания (1936), ''Воображение'' (1936) и ''Воображаемое'' (1940), а также ''Очерк теории эмоций'' (1939). Все четыре эти работы свидетельствуют о его знакомстве с философией Гуссерля. С началом Второй Мировой Войны Сартр был призван в армию и в 1940г попал в плен к немцам. После освобождения в 1941г он становится активным участником сопротивления и в 1943г публикует ''Бытие и ничто'', ясно демонстрирующий влияние Хайдеггера. В 1945г Сартр основывает литературно-политическое обозрение Les Temps Moderne (''Новые времена''), которое редактирует в течение многих лет (периодически совместно с Мерло-Понти). После войны Сартр решает прекратить преподавание и посвящает себя литературной и редакционной работе. Сартр, таким образом, один из очень немногих влиятельных философов XX века, который не работал в университете. В послевоенный период политическая активность Сартра возрастает, возрастают и его симпатии к марксизму и восхищение СССР, хотя членом коммунистической партии он так и не станет. Поддержка им СССР остается неизменной до вторжения в Венгрию в 1956г. В 1960г Сартр публикует свое второе основное произведение ''Критика диалектического разума'', что свидетельствует о его политической и социальной активности. В 1964г Сартру присудят Нобелевскую премию по литературе, но он из принципа откажется от нее. Он останется политически активным до позднего возраста, выступая против войны Франции в Алжире, активно выступая против войны во Вьетнаме вместе с Бертраном Расселом, а также поддерживая протесты в Париже в мае 1968г. На его похороны после смерти в апреле 1980г придет около 50тыс. человек.}, Морис Мерло-Понти (1908-1961)\footnote{Мерло-Понти изучал философию в престижной Высшей нормальной школе в Париже и был первым иностранным посетителем недавно созданного Архива Гуссерля в Левене в апреле 1939г. Свою первую книгу ''Структура поведения'' он публикует в 1942г, а свою, пожалуй, главную работу --- ''Феноменология восприятия'' --- в 1945г. В 1949г он становится профессором детской психологии и педагогики в Сорбонне, а в 1952г будет избран на кафедру философии в Коллеж де Франс, став самым молодым из когда-либо назначенных на эту должность, и будет занимать ее до своей смерти в мае 1961г. После Второй Мировой Войны Мерло-Понти все активнее занимается политикой и публикует ряд книг с эссе, среди которых ''Гуманизм и террор'' (1947), ''Смысл и бессмыслица'' (1948) и ''Приключения диалектики'' (1955). Параллельно с этими политическими интересами он продолжает преподавать, и многие из его лекций в Сорбонне и Коллеж де Франс свидетельствуют об обширном интересе к другим дисциплинам, включая детскую психологию, структурную лингвистику, этнологию и психоанализ. В 1960г выходит еще один том ''Знаков'', состоящий из эссе, а в 1964г будет посмертно опубликован отрывок ''Видимое и невидимое'', которое многие считают вторым главным произведением Мерло-Понти.} и Эммануэль Левинас (1906-1995)\footnote{Эммануэль Левинас родился в Каунасе в Литве (тогда входившей в состав Российской Империи). После окончания средней школы он уезжает во Францию, где начинает изучать философию в Страсбурге. В 1928г он обучается у Гуссерля и Хайдеггера во Фрайбурге, и затем совместно с Г.Пайффером переведет ''Картезианские размышления'' Гуссерля на французский. С публикацией в 1931г его докторской диссертации ''Теория интуиции в феноменологии Гуссерля'' за ним устанавливается репутация как одного из ведущих специалистов по немецкой феноменологии во Франции. К моменту начала Второй Мировой Войны Левинас стал гражданином Франции и был призван в армию. Как и Сартр, он вскоре попал в плен к немцам, где провел остаток войны, благодаря чему избежал участи большей части своей еврейской семьи, убитой немцами. После войны Левинас становится директором École Normale Israélite Orientale, а затем получает должности в университетах Пуатье (1961), Нантера (1967), а с 1973г --- в Сорбонне в Париже. Ранние послевоенные работы Левинаса ''От существования к существующему'' (1947), ''Время и другой'' (1948) и ''Открытие существования с Гуссерлем и Хайдеггером'' (1949) все еще демонстрируют его долг как перед Гуссерлем, так и перед Хайдеггером, но уже рассматривают характерные для собственного мышления Левинаса темы, такие как роль другого и отношение между этикой и онтологией. Рассмотрение этих тем достигает своей кульминации в ''Тотальность и бесконечность'' (1961) и ''Иное, чем бытие или за пределами сущности'' (1974). Помимо множества философских книг, Левинас опубликует также серию талмудических комментариев и эссе об иудаизме (в том числе ''Трудная свобода'' (1963) и ''Девять талмудических чтений'' (1968 / 1977)).}. Влияние феноменологии обусловлено в том числе возможностью понимания почти все последующих теоретических формирований в немецкой и французской философии либо как расширение феноменологии, либо как реакцию на нее. Понимание феноменологии, значит, необходимо для понимания последующего развития философии XX века.

За прошедшие годы феноменология сделала значительный вклад во многие области философии и предложила новаторский анализ интенциональности, восприятия, воплощения, эмоций, самосознания, интерсубъективности, темпоральности, историчности и истины. Она целенаправлено критикует редукционизм, объективизм и сциентизма. Подробно рассматривая человеческое существование и понимая субъекта как воплощенное и социально и культурно встроенное бытие-в-мире, феноменология делает решающий вклад в такие эмпирические дисциплины как психиатрия, социология, психология, литературоведение, антропология и архитектура.

Многие из ныне классических работ по феноменологии были написаны в первой половине XX века, но в последние годы интерес к ней возобновился. Феноменология сейчас переживает своего рода ренессанс.

Феноменология развивалась как во многом разнородное движение с множеством ответвлений, все феноменологи после Гуссерля дистанцировались от различных аспектов его исходной программы, и вряд ли корректно говорить о феноменологии как о философской системе с четко очерченным корпусом доктрин, однако не следует упускать из виду те всеобъемлющие проблемы и темы, которые объединяли и продолжают объединять ее сторонников.

В введении, конечно, невозможно отдать должное всем феноменологам, и я буду опираться прежде всего на работы Гуссерля, раннего Хайдеггера и Мерло-Понти --- трех мыслителей, чье влияние на развитие феноменологии бесспорно. Многие вводные книги по феноменологии содержат отдельные главы, посвященные этим мыслителям, но я поступлю иначе: вместо указания на их различия --- часто, на мой взгляд, преувеличивающиеся из-за фундаментально ошибочного толкования идей Гуссерля --- я сосредоточусь на том, что в них общего.

В первой части книги рассмотрена концепция философии согласно феноменологии --- вопрос метода, фокус на виде от первого лица, анализ проживаемого мира --- также развитие феноменологической традиции.

Во второй части я откладываю в сторону методологические соображения чтобы рассмотреть несколько примеров и моделей феноменологического анализа. Я рассмотрю феноменологические исследования пространственности и воплощенности, а также анализ интерсубъективности и сообщества.

Третья часть книги демонстрирует применение феноменологии вне философии. В частности, я опишу ее влияние на социологию, психологию и когнитивистику.

Главная цель этой книги --- не защита феноменологии, но ее представление: изложение ряда ее идей в максимально доступной форме без защиты их от разного рода критики. Я старался как можно меньше внимания уделять интерпретационным разногласиям и противоречиям. Так что если вы ищете систематическую защиту феноменологии и широкое соотнесение ее положений и выводов с исследовательской литературой, эта книга вам не подойдет\footnote{См., например, Захави 2003, 2005, 2014, 2017 и Галлахер и Захави 2012.}.

\part{Фундаментальные вопросы}

Следующие шесть глав представляют базовый обзор феноменологии: какого рода исследованиями она занимается, сосредоточена ли он преимущественно или даже исключительно на разуме или же в равной степени касается и мира, что такое феномен и как можно его исследовать, какие есть способы взаимодействия с миром, в чем разница между разговором о дереве и этого дерева восприятием, как мир науки соотносится со знакомым из повседневного опыта миром, что значит причастность феноменологии к трансцендентальной философии. В главах будут представлены некоторые глобальные темы и проблемы феноменологии, описаны ее методы и очерчено ее развитие.

\chapter{Феномен}

Строго говоря, феноменология занимается исследованием феноменов. Но что такое феномен, и какие феномены интересны феноменологам? Возможно, их занимают удивительные, впечатляющие, феноменальные явления? В своей автобиографии Симона де Бовуар рассказывает о первом знакомстве Сартра с феноменологией. Они были в коктейль-баре с Раймоном Ароном, только что вернувшимся из Германии, и Арон указал на заказанный им абрикосовый коктейль и сказал: ''Видите ли, мой дорогой друг, если вы феноменолог, то можете заниматься философией, говоря об этом коктейле!''. Арон был совершенно прав: даже повседневный опыт простых объектов может служить отправной точкой для феноменологического анализа. Если философия хочет избежать тупика устаревших абстракций, ей необходимо воссоединиться с богатством повседневной жизни. Впрочем, стоит отметить нацеленность феноменологии на скорее вопросы ''как'' чем ''что'' касательно объектов. Вместо того, чтобы говорить о, скажем, весе, редкости или химическом составе объекта, феноменология занимается его проявлениями, тем как он нам является. Есть существенные отличия в том, как представляют себя посуда, произведение искусства, мелодия, положение дел, число и человек. Более того, один и тот же объект может проявлять себя по-разному: с той или иной точки зрения, при сильном или слабом освещении, как воспринимаемый, воображаемый, желаемый, избегаемый, предвкушаемый или вспоминаемый. Феноменологию можно рассматривать как философский анализ различных типов данности.

\section{Будильник}

Все это звучит, наверное, абстрактно и даже немного эзотерично, так что давайте рассмотрим конкретный пример. В винтажном магазине в центре Копенгагена я ищу подарок на день рождения своему другу, и замечаю старинный будильник. Как будильник проявляет себя, что это за явление? Для начала следует признать, что простого ответа нет, поскольку будильник может проявлять себя по-разному. Я могу его не только увидеть, потрогать и услышать его, но также представить, как если бы он был дан мне на картинке, или просто-таки использовать. Для простоты давайте сосредоточимся на том, как будильник дан в рассматриваемой ситуации, т.е. как он дан в восприятии. То, как будильник будет выглядеть, зависит от освещенности --- будет это естественный солнечный свет, неоновые огни, прожекторы или что-то еще. Впрочем, вне зависимости от обстоятельств и даже при оптимальном освещении я никогда не смогу увидеть его целиком, он всегда будет дан с некой перспективы: если я посмотрю на будильник когда он стоит на столе в магазине, то смогу увидеть его верхнюю и боковые части, но не смогу увидеть заднюю, нижнюю или внутреннюю. Если же я пойду вокруг стола, то смогу увидеть обратную сторону будильника, и если подниму его со стола, то смогу осмотреть и нижнюю, но он всегда будет дан с перспективы: раскрытие одних аспектов будет сопровождаться исчезновением с поля зрения других. Это все может показаться довольно тривиальным наблюдением, повседневным и не требующим никаких дальнейших размышлений, но в руках феноменологов оно содержит ключ к множеству проблем. Прежде всего следует отметить нашу уверенность в наличие у будильника нечта большего чем то что нам о нем дано, хотя мы никогда не видим его целиком: есть его части и свойства, отсутствующие в нашем восприятии. Эти аспекты --- даже не будучи даны в восприятии --- играют в нем важную роль: воспринимая объект, мы переживаем больше, чем интуитивно данное, так что видимая сторона будильника указывает на другие, с этой перспективы не данные но могущие себя обнаружить при дальнейшем исследовании. Видимое, значит, никогда не дано изолированно, а окружено неким горизонтом, участвующем в формировании его значения.

Горизонт этот охватывает больше чем просто невидимые на мгновение аспекты рассматриваемого объекта: мы никогда не сталкиваемся с изолированными объектами, только с включенными в более широкий контекст. Так, будильник, на который я смотрю, стоит на столе в определенной комнате, в зависимости от назначения которой (торговый зал, аудитория или кабинет) будет являться по-разному, с разными значениями.

Поле моего сознания не исчерпывается будильником даже когда я наблюдаю за ним. Он может быть окружен другими часами, чашками, ручками, книгами и т.д., и сосредотачиваясь на будильнике я не буду обращать внимания на все это, но я и не буду вовсе забывать о нем, так же как и о полу, на котором стою, об освещенности комнаты и т.д.: оно будет дано как основа, контекст, части целостности фона будильника. И будучи лишь фоном они все еще могут стать объектом внимания при его смене. Возможность такого тематического изменения основана на со-данности поля основной теме так что в этом поле возможно мысленное перемещение.

Что значит данность с перспективы? В восприятии будильник дан всегда под определенным углом и на определенном расстоянии от воспринимающего. Но что это говорит нам о воспринимающем? Для такой данности будильника воспринимающий должен находиться в том же, что и будильник, пространстве, а расположенность в пространстве значит воплощенность: бестелесный воспринимающий не имел бы места в пространстве, так что будильник может быть дан только как данный воплощенному воспринимающему, и нет чисто интеллектуальной перспективы или взгляда из неоткуда --- есть только воплощенная.

Решающая роль тела в восприятии может быть проиллюстрирована и несколько иначе. Даже всегда изначально сталкиваясь с очень ограниченным взглядом на объект, мы редко остаемся таким взглядом удовлетворены. Как указывает Гуссерль, объект манит нас к дальнейшему исследованию:

\smallskip
\leftskip1cm\relax
\rightskip1cm\relax

Здесь еще есть на что посмотреть, поверни и увидь все стороны, позволь взгляду пробежаться по мне, приблизься ко мне, открой меня, раздели меня, продолжай смотреть снова и снова, поворачивай чтобы увидеть все стороны: так ты узнаешь меня, все чем я являюсь, все мои качества.

\leftskip0cm\relax
\rightskip0cm\relax
\smallskip

Как провести такое исследование, как лучше узнать будильник? Взяв его в руки и повернув, или обойдя стол так чтобы рассмотреть обратную его сторону: все это требует телесного участия и взаимодействия, предполагает их. Такое перцептивное исследование, значит --- телесная деятельность: мы двигаем глазами, головой, туловищем, руками и вообще всем телом. Эта деятельность не мгновенна, она требует времени, время играет решающую роль: если сначала посмотреть на переднюю часть будильника а затем осмотреть заднюю его часть, то передняя постепенно исчезнет из поля зрения, но все еще будет осознаваться. Такое накопление знаний о будильнике возможно именно потому что возможно сохранение виденного в прошлом. Меняя перспективу и положение, мы не воспринимаем сначала переднюю часть, затем боковую и затем заднюю, как если бы рассматривали три разных отдельных снимка. Взяв его в руки и покрутив, мы наблюдаем постепенную смену его данности, и постепенность эта значит что и наш поток сознания не серия мгновенных обособленных восприятий и имеет особую целостную временную и эмпирическую структуру. Время играет и другую роль: подчеркивая важность контекста и горизонта, о них следует думать не только в пространственных, но и во временных терминах. Мы сталкиваемся с настоящим на основе прошлого и планами и ожиданиями касательно будущего, прошлых опыт не теряется и не оставляет нас нетронутыми. Так, будильник в магазине заинтересовал меня не только потому что по опыту знаю что мой друг крепко спит, но и потому что намереваюсь сделать ему подарок.

Это подводит нас к еще одному замечанию: будильник дан мне не как приватный, не как данный лишь мне --- это публичный объект, который и другие могут наблюдать и использовать, и потому я подумал о его покупке. Даже будучи явлен мне лишь частично, он может в то же время быть явлен другим другими своими аспектами, мне недоступными.

\section{Данность и действительность}

Какое значение это все имеет? Обращая внимание на данность различных объектов, феноменология говорит лишь о явлениях, т.е. об \textit{очевидной} части природы объектов. Целью же науки обыкновенно понимается постижение объектов такими, какие они действительно есть.

В большей части философской традиции феномен определялся как представленность нам некоторого объекта, его видимость и мыслимость, и противопоставлялся объекту как он на самом деле есть. И, значит, в стремлении открыть и определить каков объект на самом деле следует выйти за рамки лишь феноменального. Понимай феноменология феномен так, она бы изучала лишь субъективное, кажущееся, \textit{поверхностное}. Как указывает Хайдеггер в седьмом разделе ''Бытия и времени'', феноменология опирается на совершенно иное, более классическое понимание феномена как того что себя проявляет, раскрывает. Феноменология, таким образом, не о лишь кажущемся, как Хайдеггер указывал также в курсе лекций, прочитанном им за несколько лет до написания ''Бытия и времени'':

\smallskip
\leftskip1cm\relax
\rightskip1cm\relax

Феноменологически абсурдно говорить о явлении как о скрывающем за собой нечто из чего оно как его представляющее и выражающее. Феномен это то за чем ничего более нет, и даже некорректно пытаться о чем-то за ним говорить потому как феноменологическая данность уже самодостаточна.

\leftskip0cm\relax
\rightskip0cm\relax
\smallskip

Некоторые могут утверждать, что феномен --- нечто субъективное, завесой дыма скрывающее объективную реальность. Феноменологи отвергают так называемую доктрину двух миров о должности принципиального различия между представленным и могущим быть понятым нами миром и миром каков он на самом деле есть. Это, конечно, не значит отрицание различия между лишь видимостью и реальностью --- в конце концов, некоторые видимости вводят в заблуждение --- но различие это не между двумя отдельными сферами, а между двумя способами явленности, между данностью при поверхностном взгляде и в результате, например, тщательного научного исследования. Феноменологи обыкновенно утверждают обладание явленным нам миром --- будь то в восприятии, повседневности или научном анализе --- всей необходимой реальностью и объективностью. Существование скрытого, закулисного и превосходящего всякую данность, мира отвергается как пустая спекуляция. Само подобное предложение, настаивают феноменологи, содержит категориальную ошибку, некорректное применение понятия реальности и злоупотребление им: вместо ее определения как недоступной и непостижимой, следует локализовать объективность внутри видимого мира.

Итак, какой же вывод следует сделать касательно области феноменологического исследования? Что поставлено на карту в феноменологическом анализе и следует ли эти исследования понимать как исследования сознания? Как теперь должно быть ясно, феноменология не фокусируется в первую очередь или даже отдельно на структуре разума: надлежащим фокусом феноменологического анализа является диада разум -- мир, или, как в конечном счете станет видно, триада я -- другой -- мир.

\chapter{Интенциональность}

Я могу думать о простых числах или далеких галактиках, представлять единорогов, чувствовать вкус ромашкового чая, бояться дементоров, вспоминать события из своего детства или узнавать датскую королеву на фотографии. Об эмпирической жизни, значит, фундаментально ошибочно думать как лишь о смеси более или менее интенсивных внутренних ощущений и чувственных состояний. Видение, слышание, воспоминание, воображение, мышление, ненависть и страх всегда о чем-то, всегда на что-то направлены --- сознание не самодавлеюще, а, скорее, самопреодолевающе. Феноменологи используют термин ''интенциональность'' для обозначения ''указания за пределы себя'', присущего сознанию.

Хотя понятие интенциональности имеет долгую историю, восходящую по крайней мере к Аристотелю, и играло центральную роль в средневековой философии, оно подвергнуто серьезному анализу в первой крупной работе по феноменологии --- ''Логические исследования'' Гуссерля (1900-1901). Гуссерль продолжал работу над этой темой до конца своей жизни, и его исследование позволило прояснить связь и различие между эмпирической субъективностью и объектами мира.

Будучи интенционально, сознание не замкнуто в себе, а занято объектами и событиями совершенно по своей природе от него отличными. Вспоминаемый мной праздничный торт, например, сильно отличается от его восприятия: он весит 3 килограмма и его можно есть, тогда как соответствующий ему опыт ничего не весит и его не получится съесть. Торт, к тому же, не о чем-то, ни на что не направлен, в отличие от восприятия.

Интенциональность аспектна и перспективна. Объект никогда не осознается как таковой, всегда определенным образом, с точки зрения или через описание. О смартфоне можно думать как о средстве связи, как о подарке, как об эффективном способе хранения музыки или как об источнике раздражения (когда он не работает должным образом). Но помимо представления может варьироваться форма представления: вместо того чтобы смотреть не таблицу, можно вообразить ее, судить о ней, запомнить ее и т.д.. Один на один и тот же объект сознание может быть направлено по-разному и одним и тем же образом оно может быть направлено на разные объекты. Например, можно сомневаться в самых разных вещах: что финансовый кризис продолжится, что выборы были честными, что климатический кризис выдуман, и можно видеть, что идет дождь, представлять что он идет, или отрицать что он идет. Каждый тип опыта --- восприятие, воображение, желание, воспоминание --- направлен на объект особым образом, и центральная задача феноменологии состоит в анализе этих различий и их взаимосвязей.

\section{Восприятие и образы}

В своих классических работах Гуссерль подчеркивает первичность перцептивной интенциональности. Например, мы могли бы обсудить, какой торт купить на день рождения: ''Отелло'' или ''Захер''? Об этих тортах мы можем думать, говорить, осознавать их даже при их в непосредственной пространственной близости отсутствии. Можем также рассмотреть фотографии торта, чтобы узнать больше о его составе и внешнем виде, при этом будучи направлены на тот же торт что и до этого, но теперь обладающий некоторым графическим присутствием. Наконец, торт мы можем купить или испечь, увидеть и попробовать. Во всех этих трех случаях мы направлены на один и тот же вид торта / вид торта, но проявляет он себя по-разному. Интенциональные акты можно ранжировать в соответствии с их способностью дать объект наиболее непосредственно. Рассмотрим, например, языковые акты: когда мы просто думаем и говорим о торте, он, конечно, объект нашего внимания, но не дан интуитивно. Изобразительные же акты уже имеют некоторое интуитивное содержание, но, подобно языковым, лишь косвенно направленны на объект: языковые акты направлены через конвенциональные представления --- языковые знаки, изобразительные --- через изображение, имеющее определенное сходство с объектом. И только акты восприятия дают нам объект непосредственно, именно этот тип направленности, по Гуссерлю, предоставляет нам сам объект в его телесном присутствии.

Утверждается, впрочем, не только представленность объекта перцептивной направленностью в привилегированном виде, но и более фундаментальный статус такой интенциональности в сравнении с интенциональностями сложными --- через воспоминание или изображение. Эпизодические воспоминания подразумевают воспоминание предыдущего опыта восприятия, воспоминание основано на восприятии. С изображениями сложнее. Стоя в Музее изящных искусств в Вене перед портретом императора Максимилиана работы Дюрера, на какой объект я направлен? Возможно, на изображение-вещь, т.е. обрамленный холст со слоями краски на нем. Возможно, на изображение-объект, т.е. на изобразительную репрезентацию и ее эстетические качества. Или, возможно, на изображение-субъект, т.е. на изображенное --- обыкновенно именно оно привлекает наше внимание первым. Может очаровать выразительный профиль Максимилиана, его роскошное платье или гранат, который он держит в левой руке --- кажется, такая работа внимания банальна и очевидна, но она включает довольно сложную стратифицированную форму интенциональности. Что значит быть изображением чего-то [другого], как один объект может изображать другой? Между знаком и обозначаемым нет внутренней связи, но предполагается наличие у изображения некоторых внутренних качеств, ограничивающих изображаемое им. При правильном сочетании буквы с, т, у, л обозначают стул, но это условность и возможно обозначение таким содержанием другого. В английском соответствующими буквами будут c, a, i, r. Но нельзя так просто решить, что на портрете Дюрера изображена банка томатного супа ''Хайнц''. Впрочем, не только наличие определенного сходства делает один предмет изображением другого: две травинки могут выглядеть совершенно одинаково, но это не делает одну из них изображением другой. Более того, сходство --- отношение двунаправленное, тогда как репрезентация однонаправлена: портрет --- изображение Максимилиана, но Максимилиан не изображение портрета. Согласно анализу изобразительной интенциональности Гуссерля, один и тот же физический объект может функционировать как изображение чего-то другого только будучи воспринимаем определенным образом. Формулируя несколько парадоксально: чтобы увидеть изображенное на картине, нужно и видеть и не видеть физически находящееся перед нами. Нужно воспринимать раму и холст, но также и выйти за пределы физически присутствующего чтобы позволить изображенному проявиться. Рамка своего рода окно в изображенное. Например, читая газету и натыкаясь на фотографию Макрона, вы фокусируете внимание на нем, но достаточно нечеткая фотография будет перетягивать направленность на себя, и даже в наилучшем случае интенциональность останется опосредованной и периферийное осознание изображения-вещи останется, иначе это будет иллюзия.

\section{Представление и причинность}

Интенциональность не просто, но многогранна и многослойна. Сознание постоянно занято объектами и событиями от него отличными, будь то через перцептивную интенциональность или через воспоминание о посещении галереи с изображением X с помощью Y. Но как это возможно? Феноменология отвергает субъективность и отсутствие влияния переживаний на внешний мир. Иначе для объяснения интенциональности был бы нужен некий интерфейс между разумом и миром: торт, на который я смотрю, осознается мной, но он, конечно, не присутствует в сознании буквально, как физический объект, а воздействует на мой сенсорный аппарат, и в результате в сознании возникает мысленное представление о торте. Т.е. восприятие предполагает две сущности: внешний сознанию объект и внутреннюю сознанию репрезентацию, и доступ к внешнему опосредован внутренним.

Такая конструкция отвергается всеми феноменологами.

Одна из их критик направлена на введение репрезентативного опосредования, вторая --- на попытку объяснить интенциональность через причинность.

Репрезентационализм, как известно, порождает скептицизм: почему это внутреннее осознание одной вещи должно способствовать осознанию совершенно другой --- внешней, и как вообще определить, представляет ли внутреннее что-то внешнее? Мы же, в конце концов, не можем выйти за пределы своего сознания и сравнить внутреннее с внешним. Кроме того, исследуя перцептивный опыт, следует понимать его как сталкивающий не с изображениями или знаками объектов (за исключением, конечно, восприятия картин или фотографий), а с самими объектами. Видя фотографию Эмпайр-Стейт-Билдинг или портрет императора Максимилиана, мы имеем дело со сложной формой интенциональности: наше осознание одного --- фотографии или картины --- позволяет осознать другое --- здание или человека, т.е. мы направлены не на воспринимаемое, а через него на что-то другое. В обычном же восприятии воспринимаемое --- не знак или изображение другого. Согласно Гуссерлю, прежде функционирования знаком или изображением объект должен быть воспринят, а значит репрезентативная теория восприятия неверна, потому что предполагает то, что пытается объяснить.

Что же касается причинности, не лучше ли объяснить перцептивную интенциональность как форму причинной связи между объектом и субъектом восприятия? Может, причинность --- клей, соединяющий разум и мир, так что сознательное состояние представляет и направлено на объект тогда и только тогда, когда связано с ним причинно? Такой довольно грубый подход сталкивается с очевидными трудностями. Во-первых, глядя в бинокль на далекий холм, я воспринимаю холм. Хотя холм (или свет, отраженный от него) мог повлиять на мою зрительную систему, это влияние было бы не единственной причиной. Почему я не воспринимаю (репрезентирую) линзы бинокля или проксимальное раздражение глазной сетчатки? Во-вторых, понятие причинности слишком грубо для схватывания аспектной природы интенциональности. Входя в комнату и видя потертый коричневый чемодан, я вижу его определенным образом, не весь сразу, в определенном освещении и на определенном фоне. Можно ли объяснить это причинным влиянием чемодана на мою зрительную систему? Следует помнить, что чемодан появляется всегда в определенном контексте и с определенным значением: в зависимости от предшествующего опыта и текущих интересов, один и тот же чемодан может проявлять себя как дорожное снаряжение, как контейнер для старых писем, как пример имеющего обратную сторону объекта, как сувенир или как символ нынешнего кризиса беженцев, и совершенно неясно, как причинное воздействие чемодана на меня способно охватить все эти различные его проявления. Но основное возражение против сведения интенциональности к причинности следующее: реально существующие пространственные объекты в моем непосредственном физическом окружении, могущие причинно на меня воздействовать --- лишь очень малая часть осознаваемого: можно думать об обратной стороне Луны, о квадратных кругах, единорогах, следующем Рождестве или принципе непротиворечия. Как эти отсутствующие, невозможные, фиктивные, будущие или идеальные объекты могли бы воздействовать на меня причинно? Возможность мышления о несуществующем кажется решающим аргументом против моего осознания объекта только лишь через его причинное на меня влияние.

Феноменологи предлагают говорить о направленности на реально существующие объекты: вкус лимона, запах кофе, холод кубика льда --- качества наличных объектов, а не каких-то ложных психических их коррелятов. Не следует говорить о перцептивном опыте как о внутреннем киноэкране, объекты не должны внутренне воспроизводиться или репрезентироваться для переживания --- перцептивный опыт дает нам объект сразу и непосредственно. Мы не переживаем репрезентации, перцептивная интенциональность не форма репрезентации --- она форма презентации.

Следует также отказаться от попытки объяснения отношения разума и мира в причинных терминах. Объекты в мире взаимодействуют причинно, но разум --- не просто объект, и его с миром отношения нельзя приравнять или ассимилировать с причинными отношениями. Подход к разуму как к такому же объекту как, например, торт, мешает раскрыть и прояснить некоторые из наиболее интересных аспектов сознания, включая его истинное эпистемическое и онтологическое значение.

Итак, феноменологи интересуются анализом структуры интенциональности для прояснения отношений между разумом и миром. Они не слишком заинтересованы отношениями между разумом и мозгом, например, и не согласны с редуктивным объяснением интенциональности через неинтенциональное. Они утверждают, что сознание характеризуется внутренней несводимой интенциональностью и что фундаментальную роль играет значение, а не причинность: камень может удариться об окно и тем самым привести к трещине в нем, эти объекты причинно связаны, но от этого не становятся одним целым. И даже если торт причинно влияет на меня, мое его перцептивное осознание основано не на причинности, а на его для меня значении.

\section{Разум и мир}

Об интенциональности можно рассказать гораздо больше, но для наших целей достаточно указания на интенциональную природу сознания и подчеркивание этим его самопревосходящего характера: разум не замкнут в ожидании причинного воздействия извне для связи с миром, и ошибочно считать мир --- чем-то внешним, а сознание --- расположенным внутри некой сферы с внутренними представлениями мира. Сознанию не нужно буквально выходить за пределы себя для достижения мира, так же как и не нужно ему этот мир как-то поглотить и переварить. Сознание не вместилище и не особое место, сознание определяется открытостью. Мы не внутри изолированной субъективной сферы, из которой время от времени выходим в мир --- интенциональная открытость сознания является его неотъемлимой частью. Это хорошо выражено в знаменитом отрывке из книги Хайдеггера ''Бытие и время'':

\smallskip
\leftskip1cm\relax
\rightskip1cm\relax

Будучи на нечто направлен и его схватывая, Dasein не будучи заключен в некой сфере выходит за пределы, но всегда уже ''вне'' вместе с этим обнаруженным в мире нечтом. Dasein не покидает никакую внутреннюю сферу, пребывая с познаваемым и определяя его характер. Даже в этом ''бытии вне'' вместе с объектом Dasein в некотором смысле ''внутри'', существует как знающее бытие-в-мире.

\leftskip0cm\relax
\rightskip0cm\relax
\smallskip

Хайдеггер не говорит об интенциональном субъекте, о субъективности или сознании, а предпочитает использовать термин ''Dasein'', состоящий из ''Da'' --- ''там'' и ''sein'' --- ''бытие''. Этот терминологический выбор частично мотивирован желанием избежать некоторых вводящих в заблуждение коннотаций, присущих традиционным терминам, и частично --- желанием подчеркнуть необходимость нахождения самого нашего существа в мире, необходимость его с ним связи. Феноменологическое исследование показывает, что традиционные категории, такие как внутреннее и внешнее, не могут адекватно отразить отношения между Dasein и миром. Поскольку Dasein всегда уже среди вещей, то не имеет внешнего и бессмысленно тогда говорить о наличии у него внутреннего. Схожие идеи можно найти и у Гуссерля, утверждающего, что поверхностное разделение между внутренним и внешним не подходит для корректного понимания интенциональности и что субъект ни в мире, ни вне него, и так же мир ни в субъекте, ни вне него. Феноменология противится самому разделению между внутренним и внешним, как выразился Мерло-Понти: ''Мир полностью внутри меня, а я полностью вне себя''.

Возвращаясь к изобразительной интенциональности, следует заметить не только ее отличие от перцептивной, но и как кажущееся естественным и простым --- видение изображения Эйфелевой башни --- в действительности довольно сложно. Мы обычно воспринимаем объекты как нечто само собой разумеющееся, не обращая внимания вклад нашей интенциональности.

\end{document}
